\begin{afterword}
\summaryhead

The post opens with an example from E.~T. Jaynes. Two researchers test the same treatment. One decided in advance to stop after 100 patients; the other decided to go on until the data showed a cure rate definitely above 60\%. Both stop with 70 cures in 100 patients. Frequentist statistics analyses the two experiments differently, because they could have ended with different data, and may find only the first result significant. The Bayesian reply is that the likelihood of the data under any hypothesis does not depend on the researcher's intentions, so the evidence and the conclusion must be the same.

The author says that Bayesians expect probability theory to be consistent and unique mathematics (Jaynes, Cox's theorem), and that any method that is not Bayesian must fail a coherence test and be open to a Dutch book. Frequentists think in tools; Bayesians think in laws. Real engines never reach the efficiency of the ideal Carnot engine, but the second law still governs them. In the same way, the Bayesian law governs every approximation, and when the exact calculation can be done, no method can do better.

\responsehead

The post's account of its opening example is right. With 100 patients fixed, 70 cures would be significant at the usual level if the true rate were 60\%; under a natural version of the second researcher's rule, they would not be. The two designs give the same likelihood ratios, so a Bayesian with a given prior draws the same conclusion from both.

The post gives the frequentist side one clause, that the experiments ``could have terminated with different data,'' and then answers a reply about tools that it writes for its opponents. The frequentist reason is that a significance level describes a procedure: how often it would give such a result if the null hypothesis were true. Run on a treatment that cures exactly 60\%, the second rule declares a rate above 60\% about 11\% of the time by patient 100 and 29\% by patient 10,000, while each of those stops, analysed as if the sample size had been fixed, would be reported at the 2.5\% level. The Bayesian side has an answer, given in Jaynes's paper: sampling to a conclusion that is appreciably false is practically impossible. At a true rate of 50\% the same rule succeeds about 1\% of the time, and no more often at 10,000 patients than at 100. The two sides disagree about which question matters: the error rate of the procedure, or the evidence in the data at hand.

That disagreement is the likelihood principle, and the post treats it as settled. Its claim that at least one of the frequentist analyses ``must discard relevant information - or simply do the wrong calculation'' follows only if the principle is granted. Birnbaum, who gave the principle its best-known derivation, later rejected it, and the derivation is still disputed.

Two general claims go beyond the results behind them. The coherence and Dutch-book arguments concern degrees of belief given to propositions; a significance test gives its hypothesis none, and the post does not show what it would fail or lose. ``You will never find a statistical method that yields a better answer'' holds on average over the calculation's own prior, which is not the frequentist standard of performance at each possible true value; the post does not say which it means. The complete class theorems, which the post does not cite, support a weaker claim: a method that no other beats everywhere is Bayesian for some prior.

In short: a correct account of the stopping-rule example and its Bayesian analysis, which gives the frequentist reasoning one clause, treats the disputed likelihood principle as settled, and states claims of coherence and optimality more broadly than the results behind them.
\end{afterword}
